\section{Harmonic geometry of the same cycles}\label{sec:cycle-metrics}

The integral identification now determines exact metric periods.
Ambient diameters require an additional argument, because a shortest
path between points of an exceptional cycle may leave it.

\subsection{Periods and primitive curve classes}
For \(z\in v^\perp\), the normalized extension gives
\[
 \theta_v(z)|_G=-(a,z)h .
\]
Thus a positive minimal line \(\ell\) represents the functional
\(-a\) on \(v^\perp\). Since \((a,v)=0\), its
Beauville--Bogomolov square is \(-2\).
The integral vector \(z_0=(0,-7f,-2)=x-2y\) satisfies
\((v,z_0)=0\), \((a,z_0)=-1\); its Mukai class restricts to
\(h\), proving primitivity of \(\ell\) directly.
The proper family \eqref{eq:line-specialization}, rather than
divisor pairings alone, identifies the moving fibre class with
this line class.

For \(t>0\), put
\begin{equation}\label{eq:kappa-scale}
 \kappa_t=\frac{88\pi t}{7},\qquad
 b_t=(0,H+tD,-88t/7).
\end{equation}
We have \((a,b_t)=-44t/7\), so
\eqref{eq:harmonic-class} gives the exact formulas
\begin{equation}\label{eq:cycle-volumes}
 \begin{gathered}
 \operatorname{Area}_{g_t}(\ell)=\operatorname{Area}_{g_t}(f)
       =\pi\kappa_t,\\
 \Vol_{g_t}(\PP^2)=\frac{(\pi\kappa_t)^2}{2},\qquad
 \Vol_{g_t}(G)=\frac{(\pi\kappa_t)^4}{12}.
 \end{gathered}
\end{equation}
Either Schubert plane type is allowed. Its \(h^2\)-integral is
one; the \(h^4\)-integral on \(G\) is two.
These are Kähler volume identities on the actual holomorphic
cycles, not a claim that the induced metric is exactly homogeneous.
For \(t<0\), set \(\kappa_t=88\pi|t|/7\) and reverse the
extension convention. The effective line functional is then
\(+a\) in that chamber's Mukai marking, and the positive
formulas are unchanged.

\subsection{A distance-preserving zero section}
Consider the positive-level complete quotient
\[
 Z_1Y_1=0,\qquad Z_1Z_1^*-Y_1^*Y_1=\kappa I_2
\]
with the flat Hilbert--Schmidt metric. Let \(g_{\Gr}\) denote
the homogeneous Grassmannian metric with minimal line area
\(\pi\).

\begin{lemma}\label{lem:grass-retraction}
The projection to \(\ker Z_1\) is a 1-Lipschitz retraction of the
quotient onto its zero section with metric \(\kappa g_{\Gr}\).
Consequently the ambient distances between core points equal
their intrinsic Grassmannian distances. In the unit normalization,
\[
 \diam G=\frac{\pi}{\sqrt2},\qquad
 \diam \PP^2=\frac{\pi}{2}
\]
for either Schubert plane, with ambient distances in the
complete quotient.
\end{lemma}

\begin{proof}
The real moment equation gives \(Z_1Z_1^*\ge\kappa I_2\).
Let \(P_{\rm row}\) be the orthogonal projector onto the row
space of \(Z_1\).
The horizontal derivative of its row-space orthonormal frame
has the form
\((I-P_{\rm row})\delta Z_1^*(Z_1Z_1^*)^{-1/2}\).
The equivalent kernel description therefore satisfies
\[
 \kappa\,|d(\ker Z_1)|_{g_{\Gr}}^2
       \le\|\delta Z_1\|^2
       \le\|\delta Z_1\|^2+\|\delta Y_1\|^2.
\]
The projection kills vertical orbit directions, so this bound
survives minimization defining the quotient norm.
On \(Y_1=0\), the equation \(Z_1Z_1^*=\kappa I_2\) gives
the Stiefel quotient metric \(\kappa g_{\Gr}\).
The retraction is the identity there; the inclusion and
retraction prove equality of the two distances.

For two planes, use principal vectors to rotate the two
orthogonal blocks through angles
\(\theta_1,\theta_2\in[0,\pi/2]\).
The horizontal frame path has length
\(\sqrt{\theta_1^2+\theta_2^2}\le\pi/\sqrt2\).
For orthogonal planes, each column of any horizontal frame
lift travels spherical distance at least \(\pi/2\).
Minkowski's integral inequality bounds its Hilbert--Schmidt
length below by \(\pi/\sqrt2\), proving sharpness.
For either Schubert plane all planes share a fixed line or
lie in a fixed three-space; one principal angle vanishes.
The same argument gives diameter \(\pi/2\) and the induced
line-area-\(\pi\) Fubini--Study metric: the principal-vector
path gives the upper bound. For the lower bound take
\(P=\langle e_1,e_2\rangle\) and
\(P'=\langle e_1,e_3\rangle\) in either Schubert plane type.
In a horizontal frame lift starting at \((e_1,e_2)\), the second
column ends in \(P'\), orthogonal to \(e_2\), and so travels
spherical distance at least \(\pi/2\). The Hilbert--Schmidt
length dominates that column length.
\end{proof}

The complete cotangent Grassmannian metric is classical
\cite[Theorem~1 and Section~2.2]{biquardgauduchon}.
The retraction argument above verifies the ambient-distance
statement needed for its transfer to the compact moduli.

\begin{proposition}\label{prop:actual-diameters}
In the actual eightfold, for either chamber,
\begin{equation}\label{eq:actual-diameters}
 \diam_{M,g_t}G=\frac{\pi}{\sqrt2}\sqrt{\kappa_t}\,
                  (1+O(|t|)),\qquad
 \diam_{M,g_t}\PP^2=\frac{\pi}{2}\sqrt{\kappa_t}\,
                  (1+O(|t|)).
\end{equation}
After distances are divided by \(\sqrt{\kappa_t}\), the ambient
spaces based on \(G\) converge in complete pointed
Gromov--Hausdorff topology to the unit-level \(T^*\Gr(2,4)\)
quotient. Along compact subsets of \(S^\circ\), the moving fibre
has transverse Eguchi--Hanson limit and
\[
 \diam_{M,g_t} f/\sqrt{\kappa_t}\longrightarrow\pi/2
\]
uniformly.
\end{proposition}

\begin{proof}
The deepest finite core has coefficient norm squared \(2\kappa_t\).
Finite competitors have length \(O(\sqrt{\kappa_t})\), and
their actual images have that length bound. The first-exit
cost from a fixed larger chart stays positive. An actual
near-minimizer cannot leave it; the 1-Lipschitz finite radius
and coarse metric bound then keep such a path in coefficient
radius \(O(\sqrt{\kappa_t})\).
On this smaller region \eqref{eq:relative} has error \(O(|t|)\).
It therefore compares both ambient distance infima, not merely
the tensors on the core. Lemma~\ref{lem:grass-retraction} gives
\eqref{eq:actual-diameters}.

After rescaling, fixed chart domains exhaust the entire finite
quotient. Exit barriers diverge, relative errors tend to zero
on every bounded model ball, and inverse coverage supplies the
whole corresponding actual ball. Finite intrinsic continuity
then gives the complete pointed limit. The unitary symmetry is
transitive on the core, so any choice of core basepoint gives
the same pointed limit up to isometry.

Fix a compact subset \(K\subset S^\circ\) and an accuracy
\(\epsilon>0\). At each point choose a full \(\C^6\) times
two-arrow quotient chart, with nested inner and outer domains,
so small that the \(\|x\|^2\) term in \eqref{eq:relative}
is below \(\epsilon\) on the outer domain. Cover \(K\) by
finitely many inner domains, as in Proposition~\ref{prop:coverage},
and then choose \(|t|\) small enough for all parameter errors
to be below \(\epsilon\) and for whole-fibre coverage to hold.
The normal zero sections are the moving projective lines,
with the same level \(\kappa_t\). Competitors of length
\(O(\sqrt{\kappa_t})\) cannot cross the positive outer exit
barriers. Projection to the normal zero section therefore gives,
uniformly over \(K\), the ratio \(\pi/2\) with a relative
error \(O(\epsilon)\). Let \(t\to0\) first and then
\(\epsilon\to0\).

The same choice works on any prescribed bounded rescaled ball.
The outer exit costs diverge after rescaling, and inverse
coverage supplies all actual points of the ball. Translating
the trivial coordinates to the moving center gives the full
pointed limit \(\C^6\) times the unit-level Eguchi--Hanson
space; its transverse factor is the asserted limit.
\end{proof}

\subsection{Successive scales with moving centers}\label{subsec:successive}
The following proposition also covers basepoints approaching a
deeper stratum.

\begin{proposition}\label{prop:successive-limits}
Assume \textup{(G1)--(G3)}.
Fix a central chart at \(p\), numbers \(\delta_n\downarrow0\),
and generic parameters \(\theta_n\to0\).
Suppose
\[
 \nu_n=\zeta_p(\theta_n)/\delta_n^2\longrightarrow\nu,
 \qquad x_n\longrightarrow x
\]
for balanced representatives \(x_n\) at level \(\nu_n\),
after unitary adjustment. Then
\begin{equation}\label{eq:first-rescaling}
 (M_{\theta_n},\delta_n^{-1}d_{\theta_n},
               f_{p,\theta_n}([\delta_nx_n]))
 \longrightarrow(Q_\nu,d_\nu,[x])
\end{equation}
in complete pointed Gromov--Hausdorff topology.

For a further scale, retain a limiting representative \(x_*\)
and its compact stabilizer \(K_*\).
Write its full normal representation as
\(W=T\oplus W'\), where \(T=W^{K_*}\) and
\(W'=T^\perp\) for the invariant hyperk\"ahler inner product.
In the common slice suppose the current representatives are
\((u_n,z_n)\to0\). For \(\sigma_n\downarrow0\), require
\begin{equation}\label{eq:successive-data}
 \frac{(\nu_n-\nu)|_{\Lie K_*}}{\sigma_n^2}\longrightarrow\lambda,
 \qquad \frac{z_n}{\sigma_n}\longrightarrow z_* .
\end{equation}
At physical scale \(\delta_n\sigma_n\), the complete pointed
limit is
\[
 T\times(W'\!/\!/\!/_\lambda K_*),\qquad
 \text{based at }(0,[z_*]).
\]
The notation \( /\!/\!/_\lambda\) means complex moment zero
and real moment \(\lambda\).
Tangential motion \(u_n\) need only tend to zero; it need not
be comparable with \(\sigma_n\).
This rule can be repeated for every fixed finite number of
stages, retaining each stabilizer, full normal representation,
normalized level and normal basepoint.
\end{proposition}

\begin{proof}
On a bounded region \(\|x\|\le R\) of the first rescaled
finite model, \eqref{eq:relative} has squared-metric error
\(O_p(|\theta_n|+\delta_n^2R^2)=o(1)\).
The fixed outer chart has a positive first-exit cost by
\eqref{eq:exit}; divided by \(\delta_n\), this tends to infinity.
Its coarse comparison and the 1-Lipschitz model radius keep
every bounded-length rescaled actual path in a bounded model
region. Conversely bounded model competitors scale into the
chart. Proposition~\ref{prop:coverage} and the inverse local
solution chart give entire actual balls. Thus both ambient
distance infima are compared. The finite intrinsic continuity
theorem \cite[Theorem~5.1]{wuebben14}, applicable by (G3), now
gives \eqref{eq:first-rescaling}.

For the next scale apply the full normal metric slice and
moving-center results \cite[Sections~4,6]{wuebben14} in this
finite model. After \((u,z)=(u_n+\sigma_n u',\sigma_n z')\),
their error on a bounded normal ball is
\[
 O_{x_*}(|\nu_n-\nu|+|u_n|+\sigma_n R)=o(1).
\]
Translation in the fixed subspace removes \(u_n\) without
changing the moment equations. The outer harmonic error is
multiplicative and is not divided by \(\sigma_n\).
Both the outer chart and a fixed inner slice have divergent
exit costs at physical scale \(\delta_n\sigma_n\).
The same two-sided path comparison and inverse coverage
therefore identify the entire complete product limit stated above.
Full normal slices retain the simple-zero hypothesis
\cite[Lemma~6.3]{wuebben14}. For any fixed finite iteration,
the relative errors sum to \(o(1)\), and the nested exit
barriers give coverage at each physical scale. This proves
the result directly from (G1)--(G3).
\end{proof}

In the eightfold, a rank-one point of the zero cone has full
normal model \(\C^6\) times the \(A_1\) quotient.
Thus approaching the deepest point is handled by the \(U(2)\)
model first, followed, when appropriate, by this full slice.
No single limit is assigned to oscillating normalized data,
and no infinite iteration is asserted.

\begin{proof}[Proof of Theorem~\ref{thm:main}]
Section~\ref{sec:construction} constructs the surface, the unique
spherical bundle, both compact chamber spaces and their common
normal target. Lemma~\ref{lem:rank-list} and the compactness
argument exclude all additional ideal boundary; the whole-fibre
and target arguments identify every stratum and fibre.
The proper family in Section~\ref{subsec:specialization}
proves reduced integral specialization.
Section~\ref{sec:schubert} gives the primitive inclusions,
self-intersection and calculated mod-six attachment.
Equations~\eqref{eq:cycle-volumes} and
\eqref{eq:actual-diameters}, with
Proposition~\ref{prop:actual-diameters}, give the cycle metric
claims. The construction verifies (G1)--(G3), so
Proposition~\ref{prop:central-completion} and
Theorem~\ref{thm:quantitative} identify the harmonic completion
and control the full compact contraction graph.
\end{proof}
